\documentclass[12pt,a4paper]{article}

% ── Encodage & langue ──────────────────────────────────────────────────────────
\usepackage[utf8]{inputenc}
\usepackage[T1]{fontenc}
\usepackage[french]{babel}

% ── Mise en page ───────────────────────────────────────────────────────────────
\usepackage[top=2.5cm, bottom=2.5cm, left=2.5cm, right=2.5cm]{geometry}
\usepackage{parskip}
\setlength{\parindent}{0pt}

% ── Typographie ────────────────────────────────────────────────────────────────
\usepackage{lmodern}
\usepackage{microtype}

% ── Couleurs & graphiques ──────────────────────────────────────────────────────
\usepackage[dvipsnames,table]{xcolor}
\usepackage{tikz}
\usetikzlibrary{patterns, calc}
\usepackage{pgfplots}
\pgfplotsset{compat=1.18}
\usepackage{pgf-pie}

% ── Tableaux ───────────────────────────────────────────────────────────────────
\usepackage{booktabs}
\usepackage{array}
\usepackage{multirow}
\usepackage{tabularx}

% ── Listes & boîtes ────────────────────────────────────────────────────────────
\usepackage{enumitem}
\usepackage{tcolorbox}
\tcbuselibrary{skins, breakable}

% ── En-têtes & pieds de page ───────────────────────────────────────────────────
\usepackage{fancyhdr}
\pagestyle{fancy}
\fancyhf{}
\fancyhead[L]{\small Projet 2}
\fancyhead[R]{\small\thepage}
\renewcommand{\headrulewidth}{0.4pt}

% ── Liens & méta ───────────────────────────────────────────────────────────────
\usepackage[hidelinks]{hyperref}
\usepackage{caption}
\usepackage{subcaption}

% ── Palettes de couleurs personnalisées ────────────────────────────────────────
\definecolor{primary}{HTML}{2C3E50}
\definecolor{accent}{HTML}{2980B9}
\definecolor{success}{HTML}{27AE60}
\definecolor{warning}{HTML}{E67E22}
\definecolor{danger}{HTML}{E74C3C}
\definecolor{lightgray}{HTML}{ECF0F1}
\definecolor{darkgray}{HTML}{7F8C8D}

% ── Commandes utilitaires ─────────────────────────────────────────────────────
\newcommand{\species}[1]{\textit{#1}}
\newcommand{\code}[1]{\texttt{\small #1}}

% ══════════════════════════════════════════════════════════════════════════════
\begin{document}

% ── Page de titre ─────────────────────────────────────────────────────────────
\begin{titlepage}
  \centering
  \vspace*{2cm}

  {\color{primary}\rule{\linewidth}{2pt}}\\[0.4cm]

  {\Huge\bfseries\color{primary} Rapport d'Analyse du Dataset}\\[0.3cm]
  {\Large\color{accent} Scraping des images d'Oiseaux e}\\[0.2cm]

  {\color{primary}\rule{\linewidth}{2pt}}\\[1.5cm]

  \begin{tcolorbox}[
    colback=lightgray, colframe=accent, arc=6pt,
    title={\bfseries Espèces étudiées}, fonttitle=\color{white},
    coltitle=white, colbacktitle=accent
  ]
    \begin{center}
      \begin{tabular}{ll}
           & — Hibou grand-duc \\[2pt]
          & — Flamant rose \\[2pt]
        & — Martin-pêcheur \\[2pt]
          & — Cygne tuberculé \\[2pt]
           & — Pic vert \\
      \end{tabular}
    \end{center}
  \end{tcolorbox}

  \vfill

  \begin{tabular}{rl}
    \textbf{Module :}    & Application IA — Semestre 6 \\
    \textbf{Projet :}    & Scrapind des images — Constitution d'un dataset d'images \\
    \textbf{Date :}      & Mai 2026 \\
  \end{tabular}

  \vspace{1cm}
  {\color{primary}\rule{\linewidth}{1pt}}
\end{titlepage}

% ── Table des matières ────────────────────────────────────────────────────────
\tableofcontents
\newpage

% ══════════════════════════════════════════════════════════════════════════════
\section{Vue d'ensemble du pipeline}

Le dataset a été constitué en quatre étapes successives :

\begin{enumerate}[label=\textbf{\arabic*.}]
  \item \textbf{Scraping} : collecte automatique d'images via Bing Images (Selenium + Requests).
  \item \textbf{Nettoyage} : suppression des images corrompues, trop petites, dupliquées ou filigranées.
  \item \textbf{Prétraitement} : redimensionnement à $224\times224$ px, conversion en JPEG, normalisation $[0,1]$.
  \item \textbf{Augmentation} : génération de $5$ variantes par image (rotation, symétrie, zoom, luminosité, contraste).
\end{enumerate}

% ══════════════════════════════════════════════════════════════════════════════
\section{Statistiques de collecte (avant nettoyage)}

\subsection{Nombre d'images collectées par espèce}

\begin{table}[h!]
  \centering
  \caption{Résultats du scraping par espèce}
  \label{tab:scraping}
  \begin{tabular}{lcccc}
    \toprule
    \textbf{Espèce} & \textbf{URLs trouvées} & \textbf{Cible} & \textbf{Collectées} & \textbf{Taux} \\
    \midrule
    Hibou grand-duc  & 105 & 100 &  87 & 87\,\% \\
    Flamant rose     & 139 & 100 & 100 & 100\,\% \\
    Martin-pêcheur   & 115 & 100 & 100 & 100\,\% \\
    Cygne tuberculé  &  70 & 100 &  68 & 68\,\% \\
    Pic vert         & 138 & 100 & 100 & 100\,\% \\
    \midrule
    \textbf{Total}   & \textbf{567} & \textbf{500} & \textbf{455} & \textbf{91\,\%} \\
    \bottomrule
  \end{tabular}
\end{table}

\subsection{Visualisation — Images collectées par espèce}

\begin{center}
\begin{tikzpicture}
  \begin{axis}[
    ybar,
    bar width=18pt,
    width=13cm, height=7cm,
    symbolic x coords={Hibou,Flamant,Martin-p.,Cygne,Pic vert},
    xtick=data,
    ymin=0, ymax=115,
    ylabel={Nombre d'images},
    xlabel={Espèce},
    nodes near coords,
    nodes near coords align={vertical},
    every node near coord/.style={font=\small},
    legend style={at={(0.98,0.98)}, anchor=north east},
    grid=major, grid style={dashed, gray!30},
    title={\textbf{Images collectées vs cible (100)}},
  ]
    \addplot[fill=accent!80, draw=accent] coordinates {
      (Hibou,87) (Flamant,100) (Martin-p.,100) (Cygne,68) (Pic vert,100)
    };
    \addplot[fill=none, draw=danger, very thick, dashed] coordinates {
      (Hibou,100) (Flamant,100) (Martin-p.,100) (Cygne,100) (Pic vert,100)
    };
    \legend{Collectées, Cible (100)}
  \end{axis}
\end{tikzpicture}
\end{center}

% ══════════════════════════════════════════════════════════════════════════════
\section{Statistiques de nettoyage}

\subsection{Tableau récapitulatif par espèce}

\begin{table}[h!]
  \centering
  \caption{Résultats du nettoyage par espèce et par motif}
  \label{tab:nettoyage}
  \renewcommand{\arraystretch}{1.3}
  \begin{tabularx}{\textwidth}{Xccccccc}
    \toprule
    \textbf{Espèce} & \textbf{Avant} & \textbf{Corrompues} & \textbf{Taille} & \textbf{Doublons} & \textbf{Filigranes} & \textbf{Total supprimé} & \textbf{Restantes} \\
    \midrule
    Hibou grand-duc  &  87 & 0 & 0 &  0 & 14 & 14 & 67\,(77\,\%) \\
    Flamant rose     & 100 & 0 & 0 &  2 &  8 & 10 & 84\,(84\,\%) \\ % 100 → 92 → 84
    Martin-pêcheur   & 100 & 0 & 0 &  0 &  5 &  5 & 93\,(93\,\%) \\
    Cygne tuberculé  &  68 & 0 & 0 &  0 & 12 & 12 & 53\,(78\,\%) \\
    Pic vert         & 100 & 0 & 0 &  1 & 20 & 21 & 78\,(78\,\%) \\
    \midrule
    \textbf{Total}   & \textbf{455} & 0 & 0 & \textbf{3} & \textbf{59} & \textbf{62} & \textbf{375} \\
    \bottomrule
  \end{tabularx}
\end{table}

\subsection{Répartition des suppressions par motif}

\begin{center}
\begin{tikzpicture}
  \pie[
    radius=3.5,
    color={danger!80, warning!80, accent!80, success!80},
    text=pin,
    before number=\phantom{},
    after number=\,\%,
    /tikz/every pin/.style={font=\small},
  ]{
    0/Corrompues (0),
    4.8/Taille insuffisante (3),
    95.2/Filigranes (59),
    0/Doublons (3)
  }
\end{tikzpicture}
\end{center}

\begin{tcolorbox}[colback=warning!10, colframe=warning, title=Observation]
  Les filigranes constituent la cause principale de suppression (\textbf{95,2\,\%} des images éliminées).
  Aucune image corrompue n'a été détectée, ce qui témoigne d'un téléchargement robuste.
\end{tcolorbox}

\subsection{Visualisation — Suppressions par espèce et motif}

\begin{center}
\begin{tikzpicture}
  \begin{axis}[
    ybar stacked,
    bar width=18pt,
    width=13cm, height=7.5cm,
    symbolic x coords={Hibou,Flamant,Martin-p.,Cygne,Pic vert},
    xtick=data,
    ymin=0, ymax=35,
    ylabel={Nombre d'images supprimées},
    legend style={at={(0.5,-0.22)}, anchor=north, legend columns=3},
    grid=major, grid style={dashed, gray!30},
    title={\textbf{Détail des suppressions par espèce}},
    nodes near coords,
    nodes near coords align={vertical},
    every node near coord/.style={font=\footnotesize},
  ]
    \addplot[fill=danger!70]   coordinates {(Hibou,0)(Flamant,0)(Martin-p.,0)(Cygne,0)(Pic vert,0)};
    \addplot[fill=warning!80]  coordinates {(Hibou,0)(Flamant,0)(Martin-p.,0)(Cygne,0)(Pic vert,0)};
    \addplot[fill=accent!70]   coordinates {(Hibou,0)(Flamant,2)(Martin-p.,0)(Cygne,0)(Pic vert,1)};
    \addplot[fill=purple!60]   coordinates {(Hibou,14)(Flamant,8)(Martin-p.,5)(Cygne,12)(Pic vert,20)};
    \legend{Corrompues, Taille, Doublons, Filigranes}
  \end{axis}
\end{tikzpicture}
\end{center}

% ══════════════════════════════════════════════════════════════════════════════
\section{Prétraitement et augmentation}

\subsection{Pipeline de prétraitement}

Chaque image nettoyée est soumise au pipeline suivant :

\begin{enumerate}[label=\textbf{\arabic*.}]
  \item \textbf{Conversion RGB} : uniformisation des canaux (suppression des modes RGBA, L, P…).
  \item \textbf{Redimensionnement} : $224 \times 224$ pixels (filtre \textsc{Lanczos}).
  \item \textbf{Normalisation} : division par 255, valeurs dans $[0, 1]$.
  \item \textbf{Sauvegarde JPEG} (qualité 95\,\%).
\end{enumerate}

\subsection{Stratégie d'augmentation}

Cinq variantes sont générées par image originale, couvrant :

\begin{itemize}
  \item Rotation aléatoire dans $[-15°,\,+15°]$
  \item Symétrie horizontale (probabilité 50\,\%)
  \item Zoom centré aléatoire ($f \in [0.9, 1.1]$)
  \item Ajustement de luminosité ($\times [0.8, 1.2]$)
  \item Ajustement de contraste ($\times [0.8, 1.2]$)
\end{itemize}

\subsection{Nombre d'images après augmentation}

\begin{table}[h!]
  \centering
  \caption{Images disponibles après prétraitement et augmentation}
  \label{tab:augmentation}
  \begin{tabular}{lcccc}
    \toprule
    \textbf{Espèce} & \textbf{Finales} & \textbf{Facteur} & \textbf{Après augmentation} \\
    \midrule
    Hibou grand-duc  &  67 & $\times 6$ &  402 \\
    Flamant rose     &  84 & $\times 6$ &  504 \\
    Martin-pêcheur   &  93 & $\times 6$ &  558 \\
    Cygne tuberculé  &  53 & $\times 6$ &  318 \\
    Pic vert         &  78 & $\times 6$ &  468 \\
    \midrule
    \textbf{Total}   & \textbf{375} & $\times 6$ & \textbf{2\,250} \\
    \bottomrule
  \end{tabular}
\end{table}

\subsection{Visualisation — Avant/après augmentation}

\begin{center}
\begin{tikzpicture}
  \begin{axis}[
    ybar,
    bar width=14pt,
    width=13cm, height=7cm,
    symbolic x coords={Hibou,Flamant,Martin-p.,Cygne,Pic vert},
    xtick=data,
    ymin=0, ymax=620,
    ylabel={Nombre d'images},
    legend style={at={(0.98,0.98)}, anchor=north east},
    grid=major, grid style={dashed, gray!30},
    nodes near coords,
    nodes near coords align={vertical},
    every node near coord/.style={font=\footnotesize},
    title={\textbf{Images finales vs images augmentées}},
  ]
    \addplot[fill=accent!80] coordinates {
      (Hibou,67)(Flamant,84)(Martin-p.,93)(Cygne,53)(Pic vert,78)
    };
    \addplot[fill=success!70] coordinates {
      (Hibou,402)(Flamant,504)(Martin-p.,558)(Cygne,318)(Pic vert,468)
    };
    \legend{Finales, Augmentées ($\times 6$)}
  \end{axis}
\end{tikzpicture}
\end{center}

% ══════════════════════════════════════════════════════════════════════════════
\section{Répartition finale Train / Validation / Test}

Une répartition classique $70\,\%$ / $15\,\%$ / $15\,\%$ est appliquée sur le jeu augmenté.

\begin{table}[h!]
  \centering
  \caption{Répartition train/val/test par espèce}
  \label{tab:split}
  \begin{tabular}{lcccc}
    \toprule
    \textbf{Espèce} & \textbf{Total} & \textbf{Train (70\,\%)} & \textbf{Val (15\,\%)} & \textbf{Test (15\,\%)} \\
    \midrule
    Hibou grand-duc  &  402 &  281 &  60 &  61 \\
    Flamant rose     &  504 &  353 &  75 &  76 \\
    Martin-pêcheur   &  558 &  391 &  83 &  84 \\
    Cygne tuberculé  &  318 &  223 &  47 &  48 \\
    Pic vert         &  468 &  328 &  70 &  70 \\
    \midrule
    \textbf{Total}   & \textbf{2\,250} & \textbf{1\,576} & \textbf{335} & \textbf{339} \\
    \bottomrule
  \end{tabular}
\end{table}

\begin{center}
\begin{tikzpicture}
  \pie[
    radius=3,
    color={accent!80, success!70, warning!80},
    text=pin,
    after number=\,\%,
  ]{70/Train (1\,576), 15/Validation (335), 15/Test (339)}
\end{tikzpicture}
\end{center}

% ══════════════════════════════════════════════════════════════════════════════
\section{Analyse critique des biais potentiels}

\subsection{Origine géographique des images}

Les images ont été collectées via \textbf{Bing Images} sans restriction géographique explicite. Cependant, plusieurs biais sont probables :

\begin{itemize}
  \item Les résultats Bing sont influencés par la langue de la requête (requêtes en \textbf{français}).
    Les images proviennent donc majoritairement de sites européens francophones (France, Belgique, Suisse).
  \item Le Hibou grand-duc et le Martin-pêcheur étant des espèces emblématiques en Europe de l'Ouest,
    leurs représentations reflètent principalement cet habitat.
  \item Le Flamant rose (présent aussi en Afrique et en Asie) pourrait être sous-représenté
    dans des contextes géographiques hors d'Europe.
\end{itemize}

\begin{tcolorbox}[colback=danger!5, colframe=danger, title=Risque identifié]
  Un modèle entraîné sur ce dataset pourrait être moins performant sur des spécimens
  photographiés hors d'Europe (sous-espèces, colorations différentes, arrière-plans exotiques).
\end{tcolorbox}

\subsection{Diversité des conditions de prise de vue}

La diversité des images Bing est \textit{a priori} plus élevée que celle d'un dataset photographié
par une seule équipe, mais certains angles morts restent probables :

\begin{itemize}
  \item \textbf{Saisons} : les oiseaux photographiés en plumage nuptial sont surreprésentés
    (plus esthétiques, donc plus partagés en ligne).
  \item \textbf{Angles} : les poses de profil ou de 3/4 dominent ; les vues en plein vol ou
    vues du dessus sont plus rares.
  \item \textbf{Arrière-plans} : les fonds naturels dominent, mais les images de volières ou
    de zoos peuvent introduire du bruit.
  \item \textbf{Éclairages} : les conditions de faible luminosité (crépuscule, nuit)
    sont quasi absentes pour les espèces diurnes.
\end{itemize}

\subsection{Disponibilité relative des espèces}

\begin{table}[h!]
  \centering
  \caption{Facilité de collecte comparée par espèce}
  \renewcommand{\arraystretch}{1.3}
  \begin{tabular}{lcp{7cm}}
    \toprule
    \textbf{Espèce} & \textbf{URLs / Collectées} & \textbf{Commentaire} \\
    \midrule
    Flamant rose    & 139 / 100 & Très iconique, abondant en ligne \\
    Pic vert        & 138 / 100 & Commun, nombreux sites naturalistes \\
    Martin-pêcheur  & 115 / 100 & Populaire en photographie animalière \\
    Hibou grand-duc & 105 / 87  & Espèce nocturne, photos rares \\
    Cygne tuberculé &  70 / 68  & Moins recherché ; beaucoup de filigranes \\
    \bottomrule
  \end{tabular}
\end{table}

\subsection{Biais de confirmation}

Le scraping par mot-clé produit un \textbf{biais de saillance} : les images les mieux classées
sont les plus représentatives (bec visible, couleurs vives, fond épuré). Les individus juvéniles,
en mue, ou partiellement cachés par la végétation sont systématiquement sous-représentés.

Ce phénomène est amplifié par :
\begin{itemize}
  \item la popularité des photos de parc/zoo sur les réseaux sociaux ;
  \item l'algorithme de Bing qui favorise les images les plus cliquées.
\end{itemize}

% ══════════════════════════════════════════════════════════════════════════════
\section{Respect des droits et aspects légaux}

\subsection{Sources des images}

Les images ont été téléchargées depuis des URLs publiquement indexées par Bing, incluant
(liste non exhaustive d'après les logs) :

\begin{itemize}
  \item \texttt{wikimedia.org} — licences Creative Commons (CC-BY, CC-BY-SA)
  \item \texttt{flickr.com} — droits variables (All Rights Reserved à CC0)
  \item \texttt{fotocommunity.com} — droits réservés (403 majoritairement rejetés)
  \item \texttt{istockphoto.com} — images sous licence commerciale (rejetées : HTTP 400/403)
  \item \texttt{animalia.bio}, \texttt{les-oiseaux.com} — droits réservés (rejetées)
  \item \texttt{facebook} (\texttt{fbsbx.com}) — droits réservés
  \item Sites de presse (\texttt{actu.fr}, \texttt{rtbf.be}) — droits réservés
\end{itemize}

\subsection{Cadre légal applicable}

\begin{tcolorbox}[colback=lightgray, colframe=primary, title=\textbf{Avertissement}]
Ce dataset est utilisé uniquement dans un cadre académique et pédagogique.  
Les images collectées restent la propriété de leurs auteurs respectifs.
\end{tcolorbox}

Quelques points légaux importants :

\begin{enumerate}[label=\textbf{\arabic*.}]
    \item \textbf{Droits d’auteur} : toute image publiée sur Internet est protégée automatiquement par le droit d’auteur.

    \item \textbf{Usage académique} : les images sont utilisées uniquement pour l’apprentissage, la recherche et l’expérimentation non commerciale.

    \item \textbf{Moteurs de recherche} : les résultats affichés par Bing ou Google ne donnent pas le droit de réutiliser librement les images.

    \item \textbf{Scraping} : la collecte automatique doit respecter les conditions d’utilisation et les règles des sites web concernés.

    \item \textbf{Protection des données} : le dataset ne contient pas volontairement d’informations personnelles sensibles.
\end{enumerate}
══════════════════════════════════════════════════════════════════════════════
\section{Synthèse}

\begin{table}[h!]
  \centering
  \caption{Tableau de synthèse global du pipeline}
  \renewcommand{\arraystretch}{1.4}
  \begin{tabular}{lc}
    \toprule
    \textbf{Étape} & \textbf{Nombre d'images} \\
    \midrule
    Collectées (scraping)              &  455 \\
    Supprimées (nettoyage)             &   62 \quad (13,6\,\%) \\
    \quad dont filigranes              &   59 \\
    \quad dont doublons                &    3 \\
    Après nettoyage                    &  375 \\
    Après prétraitement (finales)      &  375 \\
    Après augmentation ($\times 6$)    & 2\,250 \\
    \midrule
    \textbf{Jeu d'entraînement (70\,\%)} & \textbf{1\,576} \\
    \textbf{Jeu de validation (15\,\%)}  & \textbf{335} \\
    \textbf{Jeu de test (15\,\%)}        & \textbf{339} \\
    \bottomrule
  \end{tabular}
\end{table}

\vspace{1cm}

\begin{tcolorbox}[colback=success!10, colframe=success, title=\textbf{Conclusion}]
  Le dataset final contient \textbf{2\,250 images} réparties équitablement entre 5 espèces
  et 3 sous-ensembles. Le pipeline automatisé (scraping $\to$ nettoyage $\to$ prétraitement
  $\to$ augmentation) est reproductible et documenté. Les principaux biais identifiés
  (surreprésentation de l'Europe de l'Ouest, poses saillantes, plumage nuptial) devront
  être corrigés lors d'une éventuelle mise en production par l'intégration de sources
  ouvertes et diversifiées géographiquement.
\end{tcolorbox}

\end{document}